% !PRINTED 437
\weberpage{437}
Wir betrachten nun die Reihe der Klasseninvarianten
\weq{10}{\ldots, (l^{-2}k), (l^{-1}k), (k), (lk), (l^2k), (l^3k), \ldots,}
die beliebig nach vorwärts und nach rückwärts fortgesetzt werden
kann. In dieser Reihe ist, wenn $l$ zum Exponenten $\varepsilon$ gehört,
$(l^{\nu}k)$ mit $(l^{\nu+\varepsilon}k)$ identisch, während alle zwischenliegenden
Glieder davon und untereinander verschieden sind. Ist $l$ zweiseitig, so
enthält die Reihe (10) nur zwei verschiedene Glieder, von denen
jedes durch das andere \gesperrt{rational ausdrückbar} ist. Anderenfalls
schließen wir nach (8), daß jedes Glied der Reihe (10) rational
ausdrückbar ist durch die beiden vorhergehenden (oder auch durch
die beiden folgenden) in der Form
\weq{11}{(l^{\nu+1}k) = -(l^{\nu-1}k) + f_l(l^{\nu}k).}
Durch eine wiederholte Anwendung dieser Formel gelangt
man zu dem Satze, \gesperrt{daß jedes Glied der Reihe (10) rational
ausgedrückt werden kann durch irgend zwei aufeinanderfolgende
 Glieder derselben Reihe.}

\gesperrt{Ist aber $l$ zweiseitig, so sind alle Glieder der Reihe
(10) rational durch eines ausdrückbar.}

\begin{center}
§ 121. \gesperrt{Trennung der entgegengesetzten Klassen.}
\end{center}

Es kommt nun darauf an, auch im Falle eines nicht zweiseitigen
 $l$ die beiden Wurzeln der quadratischen Gleichung voneinander
 zu trennen, was durch Adjunktion von $\sqrt{D}$ möglich ist.

Ist die Klasse $l$ nicht zweiseitig, so sind die beiden Klasseninvarianten
 $(lk)$, $(l^{-1}k)$ voneinander verschieden und daher
(§ 118, III.)
\[
v = (lk) = j\left(\frac{\omega}{p}\right)
\]
eine \gesperrt{einfache} Wurzel der Transformationsgleichung
\[
F_p(u, v) = 0
\]
(§ 119, V.). Setzen wir also
\weq{1}{M = \left(\frac{\eta\left(\dfrac{\omega}{p}\right)}{\eta(\omega)}\right)^{24},}
so folgt nach § 72, daß $M$ \gesperrt{rational} durch
\[
j(\omega) \quad \text{und} \quad j\left(\frac{\omega}{p}\right),
\]
%CONT>>
